\title{LongRoPE: Extending LLM Context Window Beyond 2 Million Tokens}

\begin{abstract}
The ability to process large context windows is an important characteristic for large language models (LLMs). Nevertheless, extending the context length remains challenging because of the substantial expenses involved in fine-tuning, limited availability of sufficiently lengthy documents, and instability caused by novel token positions. As a result, state-of-the-art context window lengths remain restricted to approximately 128k tokens.

In this work, we present LongRoPE, which for the first time expands the context window of pre-trained LLMs to a remarkable \textbf{2048k} tokens. This is accomplished with up to 1k fine-tuning steps performed at training lengths of 256k, and the model continues to perform well within its original, shorter context window. Our approach is based on three core innovations: (i) we uncover and utilize two distinct types of non-uniformities in positional interpolation using an efficient search mechanism, enhancing the initial fine-tuning phase and allowing an 8$\times$ extension of context without further fine-tuning; (ii) we adopt a progressive extension technique, beginning with fine-tuning an LLM to 256k length, followed by a second positional interpolation applied to this already extended model to reach the 2048k context window; (iii) we perform a recalibration of LongRoPE at the 8k context length to restore performance for shorter context windows. Comprehensive experimental evaluations on LLaMA2 and Mistral over diverse tasks validate the efficacy of our method. Models enhanced with LongRoPE preserve their original architecture with minimal adjustment to the positional embedding, and retain compatibility with most existing optimizations. The implementation will be released at \texttt{https://github.com/microsoft/LongRoPE}

\section{Introduction}

Large Language Models (LLMs) have demonstrated impressive achievements across a range of tasks~\cite{gpt4,llama2}; however, they are constrained by the finite size of their context window, such as the 4096-token ceiling in LLaMA2~\cite{llama2}. When sequences extend beyond this window, model performance deteriorates, as these extended positions fall outside the distribution encountered during training. This constraint creates difficulties for use cases where processing extensive context is essential, such as in-context learning with many exemplars~\cite{Huang2023FewerIM} or for LLM-driven agents~\cite{park2023generative,madaan2023selfrefine}.

Recent studies have demonstrated that by fine-tuning pre-trained LLMs on longer documents, their context windows can be expanded to approximately 128k tokens~\cite{longlora,pi,yarn,zhang2024soaring,liu2023scaling}. However, pushing the boundaries of the context window faces three significant challenges. Firstly, introducing unfamiliar position indices results in numerous catastrophic values, which produce distribution shifts and hinder the fine-tuning process from converging efficiently~\cite{pi}. This problem is exacerbated in cases where the context window increases dramatically (e.g., from 4k to $>$1000k tokens), as this shift adds over 90\% entirely new position indices. Secondly, proper fine-tuning demands training texts that match the length of the extended context window. Unfortunately, most available datasets lack ample examples of extremely long texts, particularly those surpassing 1000k tokens. Furthermore, the computational cost of training on such massive sequences is very high, demanding significant time and GPU resources. Thirdly, with ultra-long context windows, the model’s attention is distributed across a vast number of tokens, which leads to diluted attention and ultimately reduces accuracy and performance when dealing with short context tasks~\cite{pi}.

To address the initial challenge, a commonly used strategy is to interpolate RoPE positional embeddings~\cite{rope,pi}, which rescales the incoming position indices so they remain within the pretrained boundaries, as illustrated in Fig.\ref{fig:overview}. The Position Interpolation (PI)~\cite{pi} method accomplishes this by applying linear interpolation to RoPE's rotary angles according to an extension factor. NTK~\cite{ntk,dynamicntk} introduces a scheme involving non-uniform interpolation and extrapolation across the RoPE embedding dimensions. YaRN~\cite{yarn} distinguishes three frequency groups within RoPE's dimensions and utilizes a combination of extrapolation, NTK-based interpolation, and linear interpolation based on these classifications. Nevertheless, within the Transformer architecture, the positional embedding possesses a \emph{complex, non-uniform information entropy}. This nuanced irregularity is not adequately harnessed in current methods, resulting in loss of information and consequently restricting the expansion of the context window.

Section~\ref{sec:analysis} empirically demonstrates two primary observations: \textbf{(1)} To achieve effective positional interpolation, it is essential to account for two specific types of non-uniformity: differences across RoPE dimensions and token positions. Reduced levels of interpolation are advantageous for both the lower RoPE dimensions and the earliest token indices; however, the best strategies vary depending on the target length to which the model is extended. 
\textbf{(2)} Incorporating these forms of non-uniformity into the positional interpolation process enables the preservation of key aspects of the original RoPE representation, notably among crucial dimensions and positions. As a result, interpolation-induced degradation is minimized, thereby providing more favorable initial conditions for subsequent fine-tuning. In addition, this approach supports an 8$\times$ sequence extension even in the absence of fine-tuning.

Building upon these results, we present \textbf{LongRoPE}, a robust approach designed to scale the LLM context window to over 2 \emph{million} tokens. LongRoPE incorporates three principal advances. Firstly, LongRoPE leverages multidimensional positional interpolation in a non-uniform manner, determining optimal rescale factors for RoPE's rotational parameters separately across each dimension, guided by token location. Given that the combinatorial space for selecting rescale factors increases exponentially as the desired extension ratio grows, LongRoPE deploys an evolutionary search strategy enhanced by two optimization techniques, substantially improving the search process. Figure~\ref{fig:overview} provides an illustration of the discovered rescaled RoPE.

Subsequently, {LongRoPE} employs a resource-efficient, incremental extension approach, enabling expansion to a 2048k context window without necessitating fine-tuning on ultra-long documents, which are rarely accessible. This method initiates by identifying suitable parameters for a 256k sequence length on the pre-trained LLM, followed by fine-tuning at this length. Given that our non-uniform positional interpolation facilitates an 8$\times$ extension in inference-only scenarios, we next determine revised RoPE rescaling parameters on this fine-tuned, length-extended model. This procedure successfully yields a 2048k context window for both LLaMA2 and Mistral~\cite{mistral}.

To address the potential drop in performance when handling the standard, smaller context windows, {LongRoPE} maintains the adaptation of the RoPE rescale parameters within the expanded LLM. As done when expanding from 256k to 2048k, we apply our search method to reduce scaling for context windows of 4k and 8k using the LLM fine-tuned at 256k, thereby aiming to minimize positional interpolation. In practice, during inference, when the input sequence is under 8k tokens, RoPE is modified in accordance with the rescale factors determined by our search process.

A broad set of experiments involving multiple LLMs and a range of tasks requiring extended context highlights the efficacy of our approach. The results indicate that LongRoPE consistently sustains low perplexity across evaluation lengths spanning from 4k to 2048k. It also attains more than 90\% accuracy in passkey retrieval, and demonstrates similar performance to existing models on conventional benchmarks structured for a 4096-token context window. The LongRoPE technique is compatible with all LLMs utilizing RoPE embeddings. We intend to publicly release both our code and LongRoPE-2048k models.

\section{Non-uniformity in Positional Interpolation}
\label{sec:analysis}

\subsection{Preliminary}

Transformer models require explicit positional information, often in the form of position embedding, to represent the order of input tokens. Our work focuses on the RoPE~\cite{rope} position embedding, which is widely used in recent  LLMs. For a token at position index $n$, its corresponding RoPE encoding can be simplified as follows:

\begin{equation}
		\fontsize{7.8}{7.8} \selectfont
	\label{eq:rope}
	\left[ cos(n\theta_0), sin(n\theta_0), cos(n\theta_1), \cdot\cdot\cdot, cos(n\theta_{d/2-1}), sin(n\theta_{d/2-1}) \right]   
\end{equation}

where $d$ denotes the dimension of the embeddings, 
 $n\theta_i$ specifies the rotary angle for the token located at position $n$, 
and $\theta_i=\theta^{-2i/d}$ encodes the rotational frequency. In the case of RoPE, the typical choice for the base value $\theta$ is 10000.

\textbf{\emph{Context window extension ratio $s$} and positional interpolation}. We introduce $s$ to represent the proportion between the new, extended context length $L'$ and the original context length $L$, defined by $s=\frac{L'}{L}$.

To extend the context window from $L$ to $L'$, current positional interpolation methods suggest downscaling rotation frequencies $\theta_i$ based on extension ratio $s$. Let $\beta=\theta^{2/d}$, and $\lambda$ denote the actual rescale factor related to $s$, we   unify these positional interpolation methods as follows:

\begin{equation}	
	\fontsize{7.5}{7.5} \selectfont
	\label{eq:unified_rope}
	\begin{aligned}
		\left[cos\left(\frac{n}{\lambda(\beta)^0}\right), sin \left(\frac{n}{\lambda(\beta)^0}\right), cos\left(\frac{n}{\lambda(\beta)^1}\right), \cdot\cdot\cdot,  sin \left(\frac{n}{\lambda(\beta)^{d/2-1}}\right) \right]   
	\end{aligned}
\end{equation}

\textbf{Linear positional interpolation (PI)}. The approach of PI~\cite{pi} implements a straightforward linear interpolation of positional indices constrained by the original pre-training sequence length. When adapting the model to a longer input via an extension ratio $s$, all rotational angles for RoPE are uniformly scaled by $\lambda=s$. Nonetheless, this uniform compression causes the positional encodings to be densely packed, making it challenging for the model to differentiate between adjacent tokens. As a result, PI's effectiveness diminishes when applied with large extension ratios.

\textbf{NTK-based interpolation and extrapolation}. ~\cite{ntk,dynamicntk} reinterpret RoPE by focusing on its role in information encoding and utilizing Neural Tangent Kernel (NTK) theory~\cite{jacot2018neural,tancik2020fourier}. In order to address the issue of densely packed positional information in PI, they advocate for spreading the interpolation demands across various RoPE dimensions. This method reduces the scaling of lower (high-frequency) dimensions, while increasing it for higher (low-frequency) dimensions, thereby enabling both positional interpolation and extrapolation scenarios where $\lambda = s^{i}$. The enhanced dynamic NTK~\cite{dynamicntk} further introduces position-dependent adjustment of the extension ratio based on the actual sequence length. In contrast to PI, which relies on fine-tuning, NTK-aware approaches are capable of expanding context windows without fine-tuning, albeit typically supporting a maximum extension ratio of 4$\times$.

\textbf{YaRN}~\cite{yarn} offers a notable enhancement in positional interpolation by partitioning RoPE dimensions into three frequency categories, each assigned a distinct interpolation method. Dimensions with high frequencies are treated using extrapolation ($\lambda$=1), low frequencies utilize linear interpolation (PI), and those in the middle range adopt the NTK technique. The central innovation of YaRN is its frequency-based dimension grouping; however, this arrangement is presently determined through manual empirical analysis, which could potentially hinder performance when applied to novel LLMs.

\subsection{Study on Non-uniform Positional Interpolation}

Drawing inspiration from NTK and YaRN, we observe that their performance improvements are largely attributable to leveraging non-linear effects—most notably, their capacity to treat frequency components along RoPE dimensions differently for targeted interpolation and extrapolation tasks. Despite these advancements, prevailing non-linear techniques are still largely constructed from hand-crafted heuristics. Consequently, this leads us to pose two critical questions: (1) Does the existing approach to positional interpolation truly represent an optimal solution? (2) Are there additional, yet-to-be-explored forms of non-linearity that could offer further benefits?

In order to address these questions, we apply evolution search (refer to Sec.~\ref{sec:method}) to identify improved non-uniform positional interpolation strategies for LLaMA2-7B. This process utilizes perplexity as the optimization criterion and is informed by evaluating 5 randomly selected instances from the PG19~\cite{pg19} validation dataset. Our experimental investigation leads to several significant observations.

\textbf{Finding 1}: \textit{RoPE dimensions exhibit substantial non-uniformities, which are not effectively handled by current positional interpolation methods.}

We determine the optimal value of $\lambda$ for each RoPE dimension as specified in Eq.~\ref{eq:unified_rope}. Table~\ref{tbl:exp1} presents a comparison of LLaMA2-7B's perplexity across various approaches on the PG19 and Proof-pile~\cite{proof-pile} test datasets, evaluated without any model fine-tuning. The results from our search yield notable performance gains, indicating that existing strategies such as the linear (PI) and non-uniform (Dynamic-NTK and YaRN) interpolation schemes do not provide optimal results. Specifically, YaRN demonstrates inferior performance relative to PI and NTK on PG19, likely because it fails to reach the intended context window length in settings where the LLM is not fine-tuned. For instance, YaRN experiences a considerable increase in perplexity beyond 7k tokens when handling an 8k context window.

In our investigation, the scaling parameters $\lambda$ in Eq.~\ref{eq:unified_rope} were observed to vary, deviating from the constant scaling factor $s$ used in PI, NTK’s formula, and YaRN’s group-specific approach. This variation in scaling leads to a substantial enhancement of LLaMA2’s language modeling capabilities, notably lowering perplexity over 8k and 16k token context lengths—even in the absence of fine-tuning. The primary reason for this improvement is that the modified positional embeddings maintain the original RoPE structure, particularly in crucial dimensions, thereby facilitating the model's ability to distinguish between adjacent token positions.

\textbf{Finding 2}: \textit{RoPE for the initial tokens in the input sequence should be extrapolated with less interpolation.} 

We propose that for the first $\hat{n}$ tokens within an input sequence, the RoPE mechanism should minimize interpolation. The rationale is that these initial tokens frequently obtain higher attention scores and thus play a pivotal role in the computation of attention, as documented in Streaming LLM~\cite{streamingllm} and LM-Infinite~\cite{han2023lminfinite}. To test this idea, we scale the context window to 8k and 16k using PI and NTK, omitting interpolation for the first $\hat n$ tokens where $\hat n$ takes values in (0, 2, ..., 256). Notably, when $\hat n$ equals 0, the configuration defaults to standard PI and NTK implementations. Table~\ref{tbl:tokenposition} reveals two significant findings: \textbf{(1)} bypassing position interpolation for the initial tokens yields a tangible boost in performance. \textbf{(2)} The choice of $\hat n$ that maximizes efficacy varies based on the desired context extension length.

\textbf{Finding 3}: \textit{Non-uniform positional interpolation effectively extends LLM context window in both fine-tuning and non-fine-tuning settings.} 

Our experiments demonstrate that the non-uniform position interpolation method we identified notably enhances extension capabilities at 8k and 16k context lengths without requiring additional training; however, even longer contexts necessitate fine-tuning. Therefore, we proceed to fine-tune LLaMA2-7B using our optimized RoPE for a 64k context window (refer to the Appendix for configuration details). Table~\ref{tbl:finetune} indicates that our approach consistently outperforms both PI and YaRN, not only on the original LLaMA2-7B but also after fine-tuning. This advantage can be attributed to our approach’s strategic non-uniform positional interpolation, which reduces information degradation and offers superior model initialization for the fine-tuning process.

\textbf{Summary}. Our findings reveal two key non-uniformities: disparities in RoPE dimensions and in token positions. Effectively exploiting these variations in positional interpolation substantially boosts the context extension performance of large language models.

\section{LongRoPE}

\label{sec:method}
Building upon these insights, we propose LongRoPE—a novel approach that initially applies an optimized search technique to leverage both forms of non-uniformity. This mechanism is subsequently utilized to scale the context window of LLMs to exceed 2 million tokens.

\subsection{Problem Formulation}

These two forms of non-uniformity significantly expand the solution space and complicate the optimization process. To manage this challenge, we recast the multidimensional non-uniform position interpolation optimization as a search-based problem.

For a LLM targeting a  context window size of $L'$ and lengthy input documents $\mathbf{X}$, where each $\mathbf{x}\in\mathbf{X}$ surpasses $L'$ in token length, we denote the original rotary angle of the $i^{th}$ dimension in RoPE embedding at token position $n$ as  $\frac{n}{\beta^i}$. The optimization problem is then formulated as follows:

\begin{equation}
\begin{gathered}
		\label{eq:problem}

\underset{\mathbf{x} \in \mathbf{X}; \, |\mathbf{x}| \geq L'}{\operatorname{arg\,min}}\, \mathcal{L}\left(\text{LLM}(\text{RoPE}, \mathbf{X})\right),\ \text{where }
\\
\scalebox{0.93}{
$\underset{\begin{subarray}{l} i = 0, \ldots, \frac{d}{2}-1; \\ n \in [0, |\mathbf{x}|); \end{subarray}}{\text{RoPE}_{(n)}} = {\left[\ldots, \cos\left(\mathbb{I}(\hat{\lambda}_{i}, \hat{n}) \cdot \frac{n}{\beta^i}\right), \sin\left(\mathbb{I}(\hat{\lambda}_{i}, \hat{n}) \cdot \frac{n}{\beta^i}\right), \ldots \right]}$}
\\
\text{where } \mathbb{I}(\hat{\lambda}_{i},\hat{n})=\begin{cases}
1 &\text{if } n < \hat{n}\\
\frac{1}{\lambda}_{i} &\text{if } n \geq \hat{n}
\end{cases}
	\end{gathered}
\end{equation}

In this approach, a set of rescale factors $\mathbb{I}(\hat{\lambda}_{i},\hat{n})$ is introduced to address two distinct types of non-uniformity. Here, $\hat{\lambda_i}$ represents dimension-specific non-uniformity within RoPE, while $\hat{n}$ serves as the threshold for positional non-uniformity among tokens. The function $\mathbb{I}(\hat{\lambda}_{i},\hat{n})$ is utilized to adjust the rotation angle associated with the $i^{th}$ RoPE dimension; $\hat\lambda_{i}$ acts as the scale modifier, and $\hat{n}$ specifies the token position threshold. For token positions preceding $\hat{n}$ (i.e., tokens $1$ through $\hat{n}-1$), the scaling factor $\hat\lambda_{i}$ is omitted, and the canonical RoPE rotational angle, given by $\frac{n}{\beta^i}$, is employed. When tokens reach positions $n\ge\hat{n}$, the rescaling factor is incorporated.

Suppose a target context window length of $L'$ is specified. The task then becomes determining the optimal set of rescaling coefficients—namely, $\mathbb{I}(\hat{\lambda}_{0},\hat{n})$, $\mathbb{I}(\hat{\lambda}_{1},\hat{n})$, ..., $\mathbb{I}(\hat{\lambda}_{i},\hat{n})$—across the $1^{st}$ through $d^{th}$ dimensions of RoPE. When these optimized scaling values are applied to the $\text{RoPE}$ in the $\text{LLM}$, the resulting model should attain the lowest possible value for the next-token prediction loss, $\mathcal{L}$ (that is, perplexity), when processing input sequences $\mathbf{X}$ with $L'$ tokens.

\subsection{Searching the Non-uniform Position Interpolation}

To address the challenge described in Eq.~\ref{eq:problem}, we propose a straightforward and robust technique that identifies optimal RoPE rescaling coefficients, allowing for comprehensive utilization of the varied, multidimensional characteristics present in position embeddings.

\textbf{Search space}. We construct an extensive search space to accommodate both sources of non-uniformity. The details of this search space are presented in Table~\ref{tbl:searchspace}. In our approach, we permit the optimization of individual rescale factors for each RoPE embedding dimension. To streamline the design, instead of directly searching for $\mathbb{I}(\hat{\lambda}_{i},\hat{n})$ (where $\hat{\lambda}_{i}=1/\lambda_i$), we focus on the parameters $\lambda_i$ and $\hat{n}$. According to Table~\ref{tbl:searchspace}, the range for $\lambda_i$ spans from a minimum of 1.0 (corresponding to simple extrapolation) up to a maximum of $s\times1.25$ (representing a broader interpolation range than PI), with increments of 0.01. Here, $s$ denotes the ratio for expanding the target context window.

The parameter $\hat{n}$ specifies how many of the starting token positions are preserved without applying position interpolation, thereby retaining their original RoPE embeddings. In practice, $\hat{n}$ is selected from the set \{0, 1, 2, 4, 8, 12, 16, 20, 24, 28, 32, 64, 128, 256\}. Setting $\hat{n}=0$ results in every token position utilizing the discovered rescale factors.

\textbf{Evolution-based search}. The search space outlined in Table~\ref{tbl:searchspace} encompasses a vast array of positional interpolation options, making systematic exploration highly complex. For instance, increasing the extension to $s=4\times$ results in $400^{128/2}\times14$ = $4\times10^{167}$ possible configurations. As the extension ratio grows, the search space proliferates at an exponential rate. To cope with this challenge, we adopt evolution search~\cite{guo2020single} and implement two optimization strategies to substantially enhance search efficiency. The complete search process is described in Algorithm~\ref{alg:search}.

\textit{Optimized initial population generation}. In place of purely random initialization for a population containing $P$ rescale factors, the initial population explicitly includes the three RoPE rescale factors linked to PI, NTK, and YaRN as distinct individuals. The rest of the initial population, amounting to $P-3$ individuals, is formed by introducing random mutations to these three rescale factors, each mutation occurring with probability $p$.

\textit{Monotonically non-decreasing constraint}. Upon generating the initial set of candidates, we evaluate each by measuring its perplexity using the target LLM, where each individual's respective RoPE rescale factors are applied to the model for processing the input $\mathbf{X}$. The best-performing $k$ candidates are then selected to serve as the parent set for subsequent evolutionary steps. Nonetheless, the sheer breadth of the search space means that unguided mutation and crossover steps may frequently yield suboptimal candidates, which can incur unnecessary computational expense from redundant perplexity evaluations. This inefficiency becomes especially prominent when $L'$ is large, given that each evaluation is resource-intensive.

To resolve this issue, we enforce a monotonic non-decreasing constraint on the sampled RoPE rescaling coefficients: $\lambda_i \leq \lambda_{i+1}$. Only those RoPE configurations that uphold this monotonicity are utilized for evaluating perplexity in the LLM, which significantly streamlines the search process. In particular, we stipulate that $\lambda_i$ must grow monotonically with respect to the RoPE dimension index ($i=0,...,63$). This requirement draws from NTK theory~\cite{jacot2018neural,tancik2020fourier,ntk}, which posits that dimensions corresponding to higher frequencies need less interpolation (and thus a lower $\lambda_i$), while those corresponding to lower frequencies allow for more interpolation (resulting in a higher $\lambda_i$).

\begin{algorithm}[t]
	\small
	\caption{The search algorithm}
	\textbf{Input:} target LLM, input samples $\mathbf{X}$, population size $P$, mutation size $N_1$, crossover size $N_2$, maximum iterations $\mathcal{T}$, mutation probability $p$\\

	\begin{algorithmic}[1]
		\label{alg:search}
		\STATE Initialize Top-k as empty set;	
		\STATE $\text{P}_0$ = \textit{Initialize\_population\_with\_optimization}($P$, $\mathbf{X}$, $p$); 
		\FOR{$i$ = $1$ to $\mathcal{T}$}
		\STATE \textit{Compute\_perplexity}(LLM, $\text{P}_{i-1}$, $\mathbf{X}$);
		\STATE Top-k $\gets$ \textit{Update\_Topk}(Top-k, $\text{P}_{i-1}$);
		\STATE $\text{P}_{mutation}$ = \textit{Mutation\_with\_mono\_constraint}(Top-k, $N_1$, $p$); 
		\STATE $\text{P}_{crossover}$ = \textit{Crossover\_with\_mono\_constraint}(Top-k, $N_2$);
		\STATE $\text{P}_i$ = $\text{P}_{mutation} \cup \text{P}_{crossover} \cup$ Top-k;
		\ENDFOR
		\STATE Output the member from Top-k exhibiting minimal perplexity;
	\end{algorithmic}
\end{algorithm}

\textbf{8$\times$ extension without fine-tuning}. Utilizing evolutionary search, we are able to determine optimal non-uniform RoPE rescale factors, maintaining essential positional and dimensional information and thereby reducing information loss caused by interpolation. As illustrated in Fig.\ref{fig:non-ft}, this approach enables LLaMA2's context window to be expanded from 4k to 32k tokens, all without the need for fine-tuning. On the other hand, methods like PI, and the non-uniform NTK and YaRN techniques, exhibit marked increases in perplexity beyond a 2$\times$ extension.

\subsection{Extending LLM Context Window to 2048K}
\label{sec:extendto2048k}

\textbf{Progressively Scaling to 2048k}. In this section, we present our approach for enlarging the context window of pre-trained LLMs from the conventional 4k up to beyond 2048k. Our experiments reveal that the proposed non-uniform positional interpolation method enables an 8$\times$ expansion without requiring fine-tuning. For even greater increases in context length (for example, 512$\times$), fine-tuning becomes essential. A common strategy involves selecting RoPE rescale factors suitable for the desired 2048k length, followed by fine-tuning. Nonetheless, this approach is hindered by the immense computational costs associated with such extensive training. Additionally, our empirical observations indicate that reliably fine-tuning LLMs for such high expansion ratios poses significant difficulties (refer to Appendix for details).

Fortunately, LongRoPE demonstrates efficacy with both pre-trained and fine-tuned extended LLMs. Consequently, we propose a progressive and efficient approach that reaches the desired 2048k length using only 1k fine-tuning iterations, while constraining training sequences to a maximum of 256k tokens.

$\Diamond$ \textit{LongRoPE search for enlarging pre-trained LLM context windows to 256k}. Using LLaMA2 as a case study, we investigate expanding the context length to both 128k and 256k tokens. The resulting expansion factors are 32$\times$ and 64$\times$ for the respective window sizes.

$\Diamond$ \textit{Fine-tuning to 256k}. Subsequently, we adapt the pre-trained LLM to support a 256k context window through fine-tuning. Initially, LLaMA2 is fine-tuned for 400 steps utilizing RoPE rescaled factors tailored for 128k. Once this stage is completed, we modify the checkpoint by updating the RoPE rescaled factors to those suitable for 256k and carry out 600 further fine-tuning steps. Compared to straightforward fine-tuning directly at 256k, this approach demonstrates greater efficiency.

$\Diamond$ \textit{Scaling fine-tuned extended LLM to 2048k using LongRoPE search}. In the last phase, we apply an additional search procedure to the previously fine-tuned LLM, which has a context length of 256k. This process expands the context window dramatically, achieving a size of 2048k while eliminating the need for additional fine-tuning steps. Consequently, the overall context length is enlarged by a factor of $512\times$.

\textbf{Mitigating degradation in the original context window}. When scaling up to a very large 2048k context window, we observe that model performance diminishes inside the initial context window range. This phenomenon is well-documented with positional interpolation~\cite{pi}, which compresses the position embeddings in the original context window into a relatively constrained area of the representation space, thereby impairing language modeling capabilities. Specifically, with a 512$\times$ extension, the embedding representations for positions across the original 4k context window are densely packed together.

To address this issue, we conduct an additional evolution search on the augmented LLM to fine-tune the RoPE rescale factors specifically for shorter context lengths (such as 4k and 8k). For these cases, we constrain the upper bound of the searched $\lambda$ values, since limited positional interpolation is needed for shorter sequences. At inference time, the LLM updates the associated RoPE rescale factors in real time.

\section{LongRoPE}

\label{sec:method}
Drawing inspiration from these results, we propose LongRoPE. This method incorporates a streamlined search algorithm designed to leverage both observed non-uniformities, subsequently enabling LLMs to achieve context windows greater than 2 million tokens.

\subsection{Problem Formulation}

These two forms of non-uniformity expand the solution space significantly and add layers of complexity to the optimization process. To tackle this issue, we reformulate the multidimensional non-uniform position interpolation optimization problem into a search problem.

For a LLM targeting a  context window size of $L'$ and lengthy input documents $\mathbf{X}$, where each $\mathbf{x}\in\mathbf{X}$ surpasses $L'$ in token length, we denote the original rotary angle of the $i^{th}$ dimension in RoPE embedding at token position $n$ as  $\frac{n}{\beta^i}$. The optimization problem is then formulated as follows:

\begin{equation}
\begin{gathered}
		\label{eq:problem}

\underset{\mathbf{x}\in\mathbf{X}; \, |\mathbf{x}|\geq L'}{\operatorname{arg\,min}} \, \mathcal{L}\, \left( \text{LLM}(\text{RoPE}, \mathbf{X})\right), \quad \text{where} \;
\\
\scalebox{0.93}{
$\underset {\begin{subarray}{l} i = 0, \ldots, \frac{d}{2} - 1;\\ n \in [0, |\mathbf{x}|); \end{subarray}}{ \text{RoPE}_ (n) } = {\left[\ldots, \cos\left(\mathbb{I}(\hat{\lambda}_i, \hat{n}) \times \frac{n}{\beta^i}\right), \sin\left(\mathbb{I}(\hat{\lambda}_i, \hat{n}) \times \frac{n}{\beta^i}\right), \ldots\right]} $}\\
\text{where} \; \mathbb{I}(\hat{\lambda}_i, \hat{n}) = \begin{cases}
	1 & \mbox{if } n < \hat{n} \\
	{\frac{1}{\lambda}_{i}} & \mbox{if } n \geq \hat{n}
\end{cases}
	\end{gathered}
\end{equation}

In this context, we define a set of scaling coefficients, $\mathbb{I}(\hat{\lambda}_{i},\hat{n})$, aimed at addressing both types of non-uniformities. Here, $\hat{\lambda}_i$ represents the irregularity across RoPE dimensions, while $\hat{n}$ indicates the irregularity with respect to token positions. The function $\mathbb{I}(\hat{\lambda}_{i},\hat{n})$ specifically adjusts the rotation angle for the $i^{th}$ dimension in the RoPE framework: $\hat{\lambda}_i$ serves as the scaling factor and $\hat{n}$ as the token position threshold. For the initial $\hat{n}$-1 tokens, the scaling factor $\hat{\lambda}_i$ is not utilized, maintaining the standard RoPE angular calculation $\frac{n}{\beta^i}$. Once token positions satisfy $n\ge\hat{n}$, the scaling factor begins to influence the computation.

Assuming a desired context window length of $L'$, the aim is to determine the set of optimal rescaling coefficients ($\mathbb{I}(\hat{\lambda}_{0},\hat{n})$, $\mathbb{I}(\hat{\lambda}_{1},\hat{n})$, ..., $\mathbb{I}(\hat{\lambda}_{i},\hat{n})$, ...) corresponding to the $1^{st}$ through $d^{th}$ dimensions of the RoPE mechanism. By applying these adjustments, the target $\text{LLM}$ utilizing the modified $\text{RoPE}$ is expected to minimize the next token prediction loss, denoted as $\mathcal{L}$ (equivalent to perplexity), on sample sequences $\mathbf{X}$ where the token length equals $L'$.

\subsection{Searching the Non-uniform Position Interpolation}

To address the challenge outlined in Eq.~\ref{eq:problem}, we present a straightforward but powerful approach that identifies optimal RoPE rescale parameters, thereby harnessing the inherent multidimensional variations within position embeddings.

\textbf{Search space}. We construct an extensive search space to capture both forms of non-uniformity. Table~\ref{tbl:searchspace} provides an overview of this space. In particular, each RoPE embedding dimension is permitted to have its own rescaling factor, which is individually optimized. To streamline the design of the search space, we opt to explore $\lambda_i$ and $\hat{n}$ as parameters, rather than directly optimizing $\mathbb{I}(\hat{\lambda}_{i},\hat{n})$, where $\hat{\lambda}_{i}=1/\lambda_i$. 
According to Table~\ref{tbl:searchspace}, the range for $\lambda_i$ starts from 1.0 (equivalent to straightforward extrapolation) up to $s\times1.25$ (corresponding to a broader interpolation window than PI), incremented by 0.01 at each step; here, $s$ denotes the context window extension ratio being targeted.

The parameter $\hat{n}$ determines how many of the earliest token positions maintain the unmodified RoPE embedding, without applying position interpolation. In our experiments, we search for optimal values of $\hat{n}$ within the set \{0, 1, 2, 4, 8, 12, 16, 20, 24, 28, 32, 64, 128, 256\}. Setting $\hat{n}=0$ indicates that every token position will utilize the rescale factors obtained through search.

\textbf{Evolution-based search}. The search space outlined in Table~\ref{tbl:searchspace} encompasses a vast array of positional interpolation configurations, making comprehensive search both computationally intensive and complex. For instance, an extension factor of $s=4\times$ yields $400^{128/2}\times14$ = 4$\times10^{167}$ potential configurations. The size of this space increases exponentially with higher extension ratios. To efficiently navigate this immense space, we employ evolution search~\cite{guo2020single} and implement two additional optimization strategies to markedly enhance search speed. Algorithm~\ref{alg:search} presents the complete search process.

\textit{Optimized initial population generation}. Rather than relying solely on random initialization for the $P$ rescale factors in the population, we explicitly include the RoPE rescale factors from PI, NTK, and YaRN as members of the initial group. The other $P$-3 individuals are created by introducing random mutations to the three rescale factors, with each factor being altered with probability $p$.

\textit{Monotonically non-decreasing constraint}. Once the initial population has been established, each member's LLM perplexity is evaluated. This involves utilizing the relevant RoPE rescale parameters within the target LLM to determine the perplexity over the input $\mathbf{X}$. The k highest-scoring individuals are then selected to serve as parents in the evolutionary process. Due to the expansive nature of the search space, indiscriminate mutation and crossover may frequently encounter suboptimal configurations, resulting in excessive perplexity evaluations. This inefficiency is magnified when $L'$ is large, since each perplexity computation requires a significant inference time.

To mitigate this issue, we enforce a non-decreasing monotonicity constraint on the sampled RoPE scaling factors, namely $\lambda_i \leq \lambda_{i+1}$. Only those RoPE configurations that adhere to this rule are used for evaluating the LLM's perplexity, thus substantially narrowing down the search space. In particular, we stipulate that $\lambda_i$ must progress monotonically as the RoPE dimension $i$ increases ($i$ = 0, ..., 63). This requirement for monotonicity across dimensions draws from NTK theory~\cite{jacot2018neural,tancik2020fourier,ntk}, which indicates that lower dimensions associated with higher frequencies demand less interpolation (i.e., smaller values of $\lambda_i$), whereas higher dimensions with lower frequencies permit greater interpolation (i.e., larger values of $\lambda_i$).

\begin{algorithm}[t]
	\small
	\caption{The search algorithm}
	\textbf{Input:} target LLM, input samples $\mathbf{X}$,   population size $P$, mutation size $N_1$, crossover size $N_2$, max iterations $\mathcal{T}$, mutate probability $p$\\

	\begin{algorithmic}[1]
		\label{alg:search}
		\STATE Top-k $\gets$ $\phi$;
		\STATE $\text{P}_0$ $\gets$ \textit{Initialize\_population\_with\_optimization} ($P$, $\mathbf{X}$, $p$);
		\FOR{$i=1$ \textbf{to} $\mathcal{T}$}
		\STATE \textit{Compute\_perplexity} (LLM, $\text{P}_{i-1}$, $\mathbf{X}$);
		\STATE Top-k $\gets$ \textit{Update\_Topk} (Top-k, $\text{P}_{i-1}$);
		\STATE $\text{P}_{mutation}$ $\gets$ \textit{Mutation\_with\_mono\_constraint} (Top-k, $N_1$, $p$);
		\STATE $\text{P}_{crossover}$ $\gets$ \textit{Crossover\_with\_mono\_constraint} (Top-k, $N_2$);
		\STATE $\text{P}_i$ $\gets$ $\text{P}_{mutation}$ $\cup$ $\text{P}_{crossover}$ $\cup$ Top-k;
		\ENDFOR
		\STATE Output the individual in Top-k possessing the minimum perplexity value;
	\end{algorithmic}
\end{algorithm}

\textbf{8$\times$ extension without fine-tuning}. Through the use of an evolutionary search approach, we successfully select non-uniform RoPE rescale factors that retain important dimensions and positions, thereby reducing the loss of information that typically results from interpolation. As shown in Fig.\ref{fig:non-ft}, our technique can increase LLaMA2’s context length from 4k to 32k without requiring any fine-tuning. By comparison, previously proposed methods—including PI and non-uniform NTK as well as YaRN—experience a sharp increase in perplexity when extended beyond 2$\times$ the original context window.

\subsection{Extending LLM Context Window to 2048K}
\label{sec:extendto2048k}

\textbf{Progressive extension to 2048k}. In this section, we present our approach for scaling the context length of pre-trained LLMs beyond the standard 4k, reaching up to 2048k tokens. Our non-uniform positional interpolation technique facilitates an extension up to 8$\times$ of the original window size without necessitating fine-tuning. However, to achieve significantly larger expansions, such as extending the window 512$\times$, fine-tuning becomes indispensable. One viable strategy involves optimizing the RoPE scaling factors specifically for the intended 2048k length, followed by subsequent fine-tuning. Nonetheless, this path is hindered by the substantial computational resources such training requires. Additionally, our empirical findings indicate that effectively fine-tuning LLMs for such large-scale extensions is particularly difficult (refer to Appendix).

LongRoPE demonstrates effectiveness on both base and fine-tuned expanded LLMs. Consequently, we propose a progressive, efficient approach that reaches the intended 2048k context length using only 1k fine-tuning steps, all conducted within a training length of 256k.

$\Diamond$ \textit{LongRoPE search applied to pre-trained LLM for 256k extension}. Using LLaMA2 as a case study, we perform a search with the goal of achieving context window lengths of 128k and 256k. This phase corresponds to extension ratios of 32$\times$ and 64$\times$, respectively.

$\Diamond$ \textit{Fine-tuning to 256k}. The next step involves adapting the pre-trained LLM to accommodate a context window of 256k tokens. This process begins with fine-tuning LLaMA2 for 400 steps utilizing RoPE rescaled factors designed for 128k. Upon completion of these steps, we substitute the RoPE rescaled factors with those for 256k on the obtained checkpoint and proceed with a further 600 steps of fine-tuning. Compared to immediate fine-tuning for 256k, this staged approach demonstrates greater efficiency.

$\Diamond$ \textit{Expanding fine-tuned extended LLM to a 2048k window using LongRoPE search}. In the last step, we apply an additional search procedure to the LLM fine-tuned with a 256k sequence length. This approach allows us to achieve a context window of 2048k, eliminating the need for additional fine-tuning. The final increase in window size amounts to a $512\times$ expansion.

\textbf{Restoring performance in shorter context windows}. Upon increasing the context window length to a sizable 2048k, we observed decreased effectiveness in the initial, smaller context window. This phenomenon aligns with previously documented limitations of positional interpolation~\cite{pi}: it compresses positional embeddings into a restricted space for the original context window, which detrimentally impacts model capability. When stretching the context by a factor of 512$\times$, the positional representations within the standard 4k window are especially densely packed.

To address this issue, we conduct an additional evolution search on the expanded LLM to fine-tune RoPE rescale factors specifically for shorter context lengths (such as 4k and 8k). For these shorter sequences, we lower the upper bound of $\lambda$ considered during the search, since less positional interpolation is necessary. At inference time, the LLM adaptively modifies the relevant RoPE rescale factors according to the input.

 \section{Experiments}

\label{sec:exp}
\subsection{Setup}
\textbf{Evaluation Tasks and Models}. LongRoPE is integrated into both LLaMA2-7B and Mistral-7B. Its effectiveness is assessed through three main criteria: (1) the perplexity scores achieved by extended-context LLMs when processing lengthy documents; (2) a Passkey retrieval task designed to test the model’s capability to extract a straightforward passkey from a large volume of irrelevant information; and (3) performance on established LLM benchmarks using a standard 4096-token context window.

\textbf{Fine-tuning}. LLaMA2 is fine-tuned utilizing a learning rate set at 2e-5, employing a linear decay schedule, and a global batch size of 32. Training proceeds for 400 steps on the Redpajama~\cite{redpajama} corpus, partitioned into 128k-length chunks marked with BOS and EOS tokens at the boundaries. Subsequently, starting from the resulting checkpoint, we conduct 600 additional steps to expand to a 256k context window. The 128k context model is trained using 8 A100 GPUs under the distributed training platform~\cite{cube}, whereas increasing to 256k context requires 16 A100 GPUs. For Mistral, training employs a fixed learning rate of 1e-6 and a global batch size of 64. Both the 128k and 256k variants use YaRN~\cite{yarn} configurations: training involves 400 steps on Together Computer's Long-Data Collections~\cite{mistraldata} with sequence lengths of 16k, using 4 A100 GPUs for these experiments.

\textbf{Search}. For target window sizes up to 256k, we set $P$=64, $N_1$=$N_2$=16, $p$=0.3, $\mathcal{T}$=40, and choose the top-32 candidates for mutation and crossover each round. Perplexity evaluation uses 5 randomly selected samples from the PG19 validation set, ensuring that each sample meets the minimum context length of the target window. For window sizes exceeding 512k, population size as well as the number of mutation and crossover candidates are reduced by half. In these cases, perplexity is computed based on 3 random validation samples from Pile-Books3~\cite{pile}.

\textbf{Baselines}. To achieve context windows of 2048k, we fine-tuned the models using 128k and 256k window sizes. Accordingly, LongRoPE-2048k (ft=128k) is applied to LLaMA2, and LongRoPE-2048k (ft=256k) is applied to Mistral. The performance of these four models is evaluated alongside leading methods for context window expansion, focusing on open-source LLMs that have undergone fine-tuning subsequent to positional interpolation via PI, NTK, and YaRN. Included in this comparison are Together-32k~\cite{together}, Code LLaMA~\cite{codellama}, LongLoRA-full-FT-100k~\cite{longlora}, YaRN-LLaMA, and YaRN-Mistral~\cite{yarn}.

\subsection{Main Results}

\label{sec:mainresult}
\textbf{Language modeling over extended sequences up to 256k tokens}. Our initial investigation involves benchmarking against leading large language models equipped with long-context capabilities, focusing on input lengths up to 256k tokens. To illustrate the robustness of our approach, we employ two distinct datasets: the test split of PG19~\cite{pile} and the Proof-pile~\cite{pg19}. Perplexity scores are assessed over varying context lengths, applying a sliding window of 256 tokens throughout the evaluation. In the case of PG19, the entire test subset consisting of 100 documents is utilized. For Proof-pile, aligning with the methodology of YaRN~\cite{yarn}, we draw 10 random samples, each exceeding a length of 128k tokens.

Table~\ref{tbl:proof-pile} and Table~\ref{tbl:pg19} present a comparison of perplexity scores for LLaMA2 and Mistral models, each enhanced through various interpolation strategies, on the Proof-pile and PG19 datasets, respectively. Two principal findings emerge: \textbf{(1)} Our extended models consistently exhibit reduced perplexity as evaluation length increases from 4k to 256k, indicating their enhanced capacity to utilize longer contexts. \textbf{(2)} Notably, even when evaluated with context windows up to 16$\times$ longer—a scenario often associated with degraded performance on shorter inputs—our LongRoPE-2048k models surpass leading baselines within the 256k context window.

\textbf{Extending language modeling to sequences longer than 2000k}. To assess performance on exceptionally lengthy texts, we utilize the Books3~\cite{pile} dataset. For practical evaluation, a random sample of 20 books, each with a length greater than 2048k, is chosen, and a sliding window approach with a span of 256k is applied.

As illustrated in Table~\ref{tbl:books3}, LongRoPE effectively expands the context window of both LLaMA2-7B and Mistral-7B models to 2048k. At shorter context lengths ranging from 8k to 128k, it delivers perplexity results that are on par with, or better than, the baselines. Differences in performance between the 2048k variants of LLaMA2 and Mistral are evident: while Mistral demonstrates superior performance relative to baselines at shorter lengths, its perplexity rises above 7 for contexts longer than 256k. On the other hand, LLaMA2 exhibits expected behavior—perplexity steadily diminishes as context length increases, although there are slight upticks observed at 1024k and 2048k. Furthermore, for LLaMA2, LongRoPE-2048k achieves improved performance when fine-tuned at a length of 256k as compared to 128k, attributed to the smaller secondary extension ratio (namely, 8$\times$ instead of 16$\times$). In contrast, Mistral yields better results with a fine-tuning window size of 128k. This discrepancy stems primarily from the application of YaRN's configuration, wherein a training length of 16k is used for both Mistral's 128k and 256k fine-tuning. This choice constrains Mistral's capability to further increase its context window following fine-tuning.

\textbf{Passkey retrieval}. Next, we investigate the practical context window size necessary for generative tasks by utilizing a synthetic evaluation protocol known as passkey retrieval, originally described by~\cite{passkey}. In this setup, the model's objective is to identify and extract a randomly placed five-digit passkey embedded somewhere within an extended document. The specific prompt format can be found in the appendix. For our experiments, we conduct 10 separate runs, with each passkey inserted at a randomly selected position, uniformly spanning the entire context window under evaluation.

Fig.~\ref{fig:passkey} presents a comparison of retrieval accuracy between our models and the baselines. For the baselines, retrieval accuracy rapidly declines to 0 once the context length surpasses 128k tokens. Conversely, although retrieving a passkey from a context containing millions of tokens is highly demanding, LongRoPE-LLaMA2-2048k (ft=256k) sustains high accuracy rates ($\ge$90\%) across the range from 4k up to 2048k tokens. Additionally, LongRoPE-Mistral-2048k (ft=128k) maintains perfect (100\%) accuracy through 1800k tokens, before decreasing to 60\% at 2048k—an outcome that corresponds with the trend observed in Table~\ref{tbl:books3}, where perplexity shows a modest rise at 2048k tokens.

\textbf{Evaluation on established benchmarks within the native context window}. The performance of LongRoPE-2048k models is assessed in their original context window using the Hugging Face Open LLM Leaderboard~\cite{huggingfaceboard}, employing both zero-shot and few-shot paradigms. Specifically, we conduct evaluations on ARC-Challenge~\cite{allenai:arc} with 25-shot, HellaSwag~\cite{zellers2019hellaswag} with 10-shot, MMLU~\cite{mmlu} with 5-shot, and TruthfulQA~\cite{lin2021truthfulqa} in the zero-shot scenario.

Table~\ref{tbl:commonsense} indicates that our models attain similar performance to previous baselines on benchmarks created for shorter context lengths, and surpass the original Mistral's score on TruthfulQA by a margin of +0.5\%. Although LongRoPE-LLaMA2-2048k, which underwent fine-tuning at 256k, experiences a minor drop in performance, its results are still acceptable across the majority of tasks.

\subsection{Ablation Results}

\textbf{Assessing the efficacy of the second positional interpolation}. Within the framework of our progressive extension methodology, we employ our search algorithm to perform a second iteration of non-uniform positional interpolation on LLMs that have already undergone fine-tuning and extension. To assess this approach, we conduct empirical evaluations using the fine-tuned LLaMA2-256k model. We further enlarge its context window to 512k, 1024k, and 2048k, comparing results for both PI and YaRN methods. The findings, presented in Table~\ref{tbl:secondsearch}, indicate that our non-uniform positional interpolation successfully preserves perplexity at a stable rate as the context window extends. In comparison, both PI and YaRN experience a marked increase in perplexity as the extension ratio grows.

\textbf{Enhancing performance at reduced context lengths.} In order to address performance degradation at shorter contexts, we optimize the RoPE scaling factors for LongRoPE-2048k using our dedicated search method. We specifically restrict the maximum scale factors considered during the search, thereby limiting interpolation when applied to 4k and 8k contexts. Table~\ref{tbl:shortperformance} presents perplexity data for LongRoPE-LLaMA2-2048k evaluated on Proof-pile at 4k and 8k, as well as the average accuracy from LLM benchmarks. These results indicate a pronounced improvement in effectiveness for shorter context windows.

\textbf{Investigation of both types of non-uniformity}. In this section, we examine the individual effect of the two identified non-uniformities through ablation studies. Two experimental setups are considered: (i) we scale LLaMA2-7B up to relatively short sequence lengths (16k and 32k) employing various techniques—position interpolation (PI), optimizing solely the RoPE dimension, and optimizing both non-uniformities; (ii) we further increase the length of our fine-tuned 256k LLaMA2 model to 2048k using the same approach. Perplexity measurements are recorded without additional fine-tuning. According to Table~\ref{tbl:searchdimension}, adjusting non-uniformity in the RoPE dimension leads to a notable decrease in perplexity compared to simple linear interpolation with PI. Moreover, non-uniformity in token position yields noticeable gains in model performance at 16k and 32k sequence lengths; however, its influence diminishes at 2048k, likely due to the exceptionally lengthy input. In such cases, preserving only the earliest tokens without interpolation appears to lose effectiveness, a subject we intend to explore further in subsequent research.

\section{Related Works}

Besides position interpolation methods, this section also examines relevant literature employing alternative strategies.

\textbf{Retrieval-based} techniques incorporate an external memory mechanism to store extensive historical context and utilize retrieval components to obtain pertinent documents during inference~\cite{longllama,wang2023augmenting,borgeaud2022improving}. These approaches generally require substantial changes to the architecture of LLMs. In comparison, our method is more efficient and minimally invasive, only necessitating slight adjustments to positional embeddings. Furthermore, our approach extends to a broader set of long-context applications beyond mere retrieval, encompassing tasks such as lengthy document summarization and few-shot learning.

\textbf{Attention-based context window extensions}. In addition to interpolating positional embeddings, alternative approaches expand the input context capacity while maintaining the LLM’s original context window by adapting the underlying attention mechanism~\cite{han2023lminfinite, streamingllm,parallelwindow}. Central to these techniques is the use of specialized attention masks to address the challenge of increased attention computation due to novel positions. Such strategies can be used in conjunction with positional interpolation methods, as they are mutually supportive.

\textbf{Fine-tuning based approaches} aim to enhance the fine-tuning process of pre-trained LLMs by adapting their position embeddings for extended context lengths. Approaches such as Code LLaMA~\cite{codellama}, LLaMA2 Long~\cite{xiong2023effective}, and ScaledRoPE~\cite{liu2023scaling} utilize an enlarged base value for RoPE and subsequently fine-tune at the intended sequence length. In contrast, our approach allows for greater adaptability across a range of target sequence lengths and is capable of supporting sequences exceeding 2M tokens. Given the significant GPU requirements of fine-tuning for extremely long contexts (e.g., beyond 128k tokens), more recent methods like LongLoRA~\cite{longlora} and PoSE~\cite{zhu2023pose} have been introduced to alleviate such computational challenges. Our technique is complementary to these efficient fine-tuning strategies.

\section{Conclusion}

This paper introduces LongRoPE, a technique that substantially increases the context window of LLMs up to a groundbreaking 2048k tokens, all while preserving their proficiency within conventional short-context settings. Our approach utilizes two distinct non-uniformities found in RoPE positional embeddings, which we efficiently discover through evolutionary search. This methodology yields dual advantages: it not only furnishes an advantageous starting point for subsequent fine-tuning but also supports a context window expansion by a factor of 8$\times$ without requiring any fine-tuning. Furthermore, we propose a gradual extension strategy that leverages LLMs fine-tuned at 256k token lengths to ultimately achieve a 2048k token window, circumventing additional fine-tuning steps. Comprehensive experimental results demonstrate the robustness and efficacy of LongRoPE. We anticipate that LongRoPE-2048k models will unlock new possibilities for long-context applications and promote future advancements in this area.

\section*{Broader Impacts} The purpose of this research is to contribute to the progress of Machine Learning as a discipline. While our findings may give rise to various societal effects, we do not identify any particular outcomes that warrant special attention within this section.

	\balance

\end{document}